\documentclass[10pt]{extarticle}
\usepackage[a4paper,margin=1.5cm]{geometry}
\usepackage[T1]{fontenc}
\usepackage{amsmath,amssymb}
\usepackage{graphicx}
\usepackage[export]{adjustbox}
\usepackage{multicol}
\usepackage{enumitem}
\usepackage{booktabs,array,tabularx}
\setlength{\columnsep}{0.7cm}
\setlength{\parindent}{0pt}
\setlength{\parskip}{3pt}
\pagestyle{empty}
\makeatletter
% figure registry (figdefs.tex is generated by the extract step)
\newcommand{\qfigdef}[4]{\@namedef{qf@file@#1}{#2}\@namedef{qf@w@#1}{#3}\@namedef{qf@wide@#1}{#4}}
\newcommand{\qf@check}[1]{\@ifundefined{qf@file@#1}{\PackageError{score2}{Unknown figure id #1}{}}{}}
\newcommand{\qf@img}[4][]{\includegraphics[width=#2\dimexpr\@nameuse{qf@w@#4}pt\relax,max width=#3,max height=0.45\textheight,#1]{\@nameuse{qf@file@#4}}}
\newif\ifqf@full
\newif\ifqf@boxed
\newsavebox{\qf@box}
\newcommand{\qf@open}{\global\qf@fulltrue\par\global\qf@boxedtrue
  \setbox\qf@box\vbox\bgroup\hsize\textwidth\linewidth\textwidth}
\newcommand{\qf@emit}[1]{\par\egroup\global\qf@boxedfalse
  \ifdim\dimexpr\ht\qf@box+\dp\qf@box\relax<#1\textheight\noindent\usebox{\qf@box}\par\else\unvbox\qf@box\fi}
\newcommand{\qfullsplit}{\ifqf@boxed\qf@emit{0.85}\fi}
\newcommand{\qf@close}{\ifqf@boxed\qf@emit{0.45}\fi\global\qf@fullfalse}
\newcommand{\qfullbegin}{\end{multicols}\qf@open}
\newcommand{\qfullend}{\qf@close\begin{multicols}{2}}
\newcommand{\qfullstart}{\qf@open}
\newcommand{\qfulldone}{\qf@close}
\newcommand{\qfig}[2][1]{\qf@check{#2}%
  \ifqf@full\qf@blk{#1}{#2}\else
  \ifnum\@nameuse{qf@wide@#2}=1 \expandafter\@firstoftwo\else\expandafter\@secondoftwo\fi
  {\end{multicols}\begin{center}\qf@img{#1}{\textwidth}{#2}\end{center}\begin{multicols}{2}}%
  {\qf@blk{#1}{#2}}\fi}
\newcommand{\qf@blk}[2]{\par\smallskip{\centering\qf@img{#1}{\linewidth}{#2}\par}\smallskip}
\newcommand{\qfigi}[2][1]{\qf@check{#2}\qf@img[valign=c]{#1}{\linewidth}{#2}}
\makeatother
% structure macros - lint reads these to check completeness and figure placement
\newcommand{\papersec}[1]{\par\bigskip{\centering\Large\bfseries PAPER-#1\par}\medskip}
\newcommand{\partsec}[2]{\par\medskip\noindent{\large\bfseries PART-#1: #2}\par\smallskip}
\newcommand{\sect}[1]{\par\smallskip\noindent\textbf{#1}\par}
\newcommand{\qn}[1]{\par\smallskip\noindent\textbf{#1.}~}
\newcommand{\paraq}[2]{\par\smallskip\noindent\textbf{Paragraph for Questions #1 and #2}\par}
\newenvironment{qoptions}{\begin{itemize}[leftmargin=1.8em,itemsep=1pt,topsep=2pt,label={}]}{\end{itemize}}
\newcommand{\opt}[2]{\item[(#1)] #2}
% figure registry (generated by score2_tools.py finalize)
\qfigdef{p1_phy_s1_q04_1}{Score-II - Test 13 (06 Feb 2024)-img/p1_phy_s1_q04_1.png}{94.8}{0}
\qfigdef{p1_phy_s1_q05_1}{Score-II - Test 13 (06 Feb 2024)-img/p1_phy_s1_q05_1.png}{56.8}{0}
\qfigdef{p1_phy_s1_q05_2}{Score-II - Test 13 (06 Feb 2024)-img/p1_phy_s1_q05_2.png}{56.8}{0}
\qfigdef{p1_phy_s1_q05_3}{Score-II - Test 13 (06 Feb 2024)-img/p1_phy_s1_q05_3.png}{56.8}{0}
\qfigdef{p1_phy_s1_q05_4}{Score-II - Test 13 (06 Feb 2024)-img/p1_phy_s1_q05_4.png}{56.8}{0}
\qfigdef{p1_phy_s1_q06_1}{Score-II - Test 13 (06 Feb 2024)-img/p1_phy_s1_q06_1.png}{120.5}{0}
\qfigdef{p1_phy_s1_q07_1}{Score-II - Test 13 (06 Feb 2024)-img/p1_phy_s1_q07_1.png}{76.3}{0}
\qfigdef{p1_phy_s1_q08_1}{Score-II - Test 13 (06 Feb 2024)-img/p1_phy_s1_q08_1.png}{95.0}{0}
\qfigdef{p1_phy_s1_q10_1}{Score-II - Test 13 (06 Feb 2024)-img/p1_phy_s1_q10_1.png}{79.8}{0}
\qfigdef{p1_phy_s1_q13_1}{Score-II - Test 13 (06 Feb 2024)-img/p1_phy_s1_q13_1.png}{78.6}{0}
\qfigdef{p1_phy_s1_q19_1}{Score-II - Test 13 (06 Feb 2024)-img/p1_phy_s1_q19_1.png}{136.3}{0}
\qfigdef{p1_phy_s2_q02_1}{Score-II - Test 13 (06 Feb 2024)-img/p1_phy_s2_q02_1.png}{110.0}{0}
\qfigdef{p1_phy_s2_q03_1}{Score-II - Test 13 (06 Feb 2024)-img/p1_phy_s2_q03_1.png}{97.4}{0}
\qfigdef{p1_phy_s2_q05_1}{Score-II - Test 13 (06 Feb 2024)-img/p1_phy_s2_q05_1.png}{90.6}{0}
\qfigdef{p1_chem_s1_q15_1}{Score-II - Test 13 (06 Feb 2024)-img/p1_chem_s1_q15_1.png}{78.7}{0}
\qfigdef{p1_chem_s1_q15_2}{Score-II - Test 13 (06 Feb 2024)-img/p1_chem_s1_q15_2.png}{34.9}{0}
\qfigdef{p1_chem_s1_q15_3}{Score-II - Test 13 (06 Feb 2024)-img/p1_chem_s1_q15_3.png}{43.0}{0}
\qfigdef{p1_chem_s1_q15_5}{Score-II - Test 13 (06 Feb 2024)-img/p1_chem_s1_q15_5.png}{50.1}{0}
\qfigdef{p1_chem_s1_q15_4}{Score-II - Test 13 (06 Feb 2024)-img/p1_chem_s1_q15_4.png}{37.5}{0}
\qfigdef{p1_chem_s1_q16_1}{Score-II - Test 13 (06 Feb 2024)-img/p1_chem_s1_q16_1.png}{110.9}{0}
\qfigdef{p1_chem_s1_q16_2}{Score-II - Test 13 (06 Feb 2024)-img/p1_chem_s1_q16_2.png}{48.8}{0}
\qfigdef{p1_chem_s1_q16_3}{Score-II - Test 13 (06 Feb 2024)-img/p1_chem_s1_q16_3.png}{47.3}{0}
\qfigdef{p1_chem_s1_q16_5}{Score-II - Test 13 (06 Feb 2024)-img/p1_chem_s1_q16_5.png}{48.8}{0}
\qfigdef{p1_chem_s1_q16_4}{Score-II - Test 13 (06 Feb 2024)-img/p1_chem_s1_q16_4.png}{55.1}{0}
\qfigdef{p1_chem_s1_q17_2}{Score-II - Test 13 (06 Feb 2024)-img/p1_chem_s1_q17_2.png}{76.5}{0}
\qfigdef{p1_chem_s1_q17_1}{Score-II - Test 13 (06 Feb 2024)-img/p1_chem_s1_q17_1.png}{49.2}{0}
\qfigdef{p1_chem_s1_q17_3}{Score-II - Test 13 (06 Feb 2024)-img/p1_chem_s1_q17_3.png}{54.0}{0}
\qfigdef{p1_chem_s1_q17_4}{Score-II - Test 13 (06 Feb 2024)-img/p1_chem_s1_q17_4.png}{45.6}{0}
\qfigdef{p1_chem_s1_q18_1}{Score-II - Test 13 (06 Feb 2024)-img/p1_chem_s1_q18_1.png}{161.4}{0}
\qfigdef{p1_chem_s1_q18_2}{Score-II - Test 13 (06 Feb 2024)-img/p1_chem_s1_q18_2.png}{31.1}{0}
\qfigdef{p1_chem_s1_q18_3}{Score-II - Test 13 (06 Feb 2024)-img/p1_chem_s1_q18_3.png}{33.2}{0}
\qfigdef{p1_chem_s1_q18_4}{Score-II - Test 13 (06 Feb 2024)-img/p1_chem_s1_q18_4.png}{26.3}{0}
\qfigdef{p1_chem_s1_q18_5}{Score-II - Test 13 (06 Feb 2024)-img/p1_chem_s1_q18_5.png}{55.4}{0}
\qfigdef{p1_chem_s1_q19_1}{Score-II - Test 13 (06 Feb 2024)-img/p1_chem_s1_q19_1.png}{132.4}{0}
\qfigdef{p1_chem_s1_q20_1}{Score-II - Test 13 (06 Feb 2024)-img/p1_chem_s1_q20_1.png}{157.6}{0}
\qfigdef{p1_chem_s1_q20_A}{Score-II - Test 13 (06 Feb 2024)-img/p1_chem_s1_q20_A.png}{55.3}{0}
\qfigdef{p1_chem_s1_q20_3}{Score-II - Test 13 (06 Feb 2024)-img/p1_chem_s1_q20_3.png}{54.2}{0}
\qfigdef{p1_chem_s1_q20_C}{Score-II - Test 13 (06 Feb 2024)-img/p1_chem_s1_q20_C.png}{50.1}{0}
\qfigdef{p1_chem_s1_q20_5}{Score-II - Test 13 (06 Feb 2024)-img/p1_chem_s1_q20_5.png}{75.5}{0}
\qfigdef{p1_chem_s2_q04_1}{Score-II - Test 13 (06 Feb 2024)-img/p1_chem_s2_q04_1.png}{96.8}{0}
\qfigdef{p1_chem_s2_q08_1}{Score-II - Test 13 (06 Feb 2024)-img/p1_chem_s2_q08_1.png}{51.3}{0}
\qfigdef{p1_chem_s2_q09_1}{Score-II - Test 13 (06 Feb 2024)-img/p1_chem_s2_q09_1.png}{81.5}{0}
\qfigdef{p1_chem_s2_q10_1}{Score-II - Test 13 (06 Feb 2024)-img/p1_chem_s2_q10_1.png}{126.5}{0}
\begin{document}
\begin{center}{\Large\bfseries JEE (Main+Advanced) Score-II --- Test 13 (06 Feb 2024)}\end{center}
\begin{multicols}{2}
\partsec{1}{PHYSICS}

\sect{Section-I: Single Correct}

\qn{1} Mass defect in a $_{11}Na^{23}$ atom is 0.21 amu. Then the binding energy per nucleon will be :-
\begin{qoptions}
\opt{A}{8.505 MeV}
\opt{B}{8.605 MeV}
\opt{C}{9.23 MeV}
\opt{D}{9.315 MeV}
\end{qoptions}

\qn{2} When a monochromatic point source of light is at a distance of 0.2 m from a photo-electric cell, the stopping potential and the saturation current are respectively 0.6 volt and 18.0 mA. If the same source is placed 0.6 m away from the photoelectric cell, then :-\\
(a) The stopping potential will be 0.2 volt\\
(b) The stopping potential will be 0.6 volt\\
(c) The saturation current will be 6.0 mA\\
(d) The saturation current will be 2.0 mA
\begin{qoptions}
\opt{A}{a, c}
\opt{B}{a, d}
\opt{C}{b, c}
\opt{D}{b, d}
\end{qoptions}

\qn{3} The energy of a photon of wavelength 6600 \AA\ in watt-hour is :-
\begin{qoptions}
\opt{A}{$3\times10^{-19}$}
\opt{B}{$1.875\times10^{-25}$}
\opt{C}{$8.33\times10^{-23}$}
\opt{D}{$5\times10^{-21}$}
\end{qoptions}

\qn{4} Current drawn from the cell in the given circuit :-
\qfig{p1_phy_s1_q04_1}
\begin{qoptions}
\opt{A}{$\dfrac16$ A}
\opt{B}{$\dfrac12$ A}
\opt{C}{$\dfrac13$ A}
\opt{D}{2A}
\end{qoptions}

\qn{5} A point charge q is placed at a point inside a hollow conducting sphere. Which of the following electric force pattern is correct?
\begin{qoptions}
\opt{A}{\qfigi{p1_phy_s1_q05_1}}
\opt{B}{\qfigi{p1_phy_s1_q05_2}}
\opt{C}{\qfigi{p1_phy_s1_q05_3}}
\opt{D}{\qfigi{p1_phy_s1_q05_4}}
\end{qoptions}

\qn{6} Shown in the figure is a circular loop of radius r and resistance R. A variable magnetic field of induction $B=B_0e^{-t}$ is established inside the coil. If the key (K) is closed, the electrical power developed right after closing the switch is equal to :-
\qfig{p1_phy_s1_q06_1}
\begin{qoptions}
\opt{A}{$\dfrac{B_0^2\pi r^2}{R}$}
\opt{B}{$\dfrac{B_010r^3}{R}$}
\opt{C}{$\dfrac{B_0^2\pi^2r^4R}{5}$}
\opt{D}{$\dfrac{B_0^2\pi^2r^4}{R}$}
\end{qoptions}

\qn{7} For the circuit shown in the figure, the current through the inductor is 0.9 A while the current through the condenser is 0.4A. Hence the current drawn from the generator is:
\qfig{p1_phy_s1_q07_1}
\begin{qoptions}
\opt{A}{I = 1.13 A}
\opt{B}{I = 0.9 A}
\opt{C}{I = 0.5 A}
\opt{D}{I = 0.6 A}
\end{qoptions}

\qn{8} In the circuit shown in figure reading of voltmeter is $V_1$ when only $K_1$ is closed, reading of voltmeter is $V_2$ when only $K_2$ is closed and reading of voltmeter is $V_3$ when both $K_1$ and $K_2$ are closed. Then :-
\qfig{p1_phy_s1_q08_1}
\begin{qoptions}
\opt{A}{$V_3>V_2>V_1$}
\opt{B}{$V_2>V_1>V_3$}
\opt{C}{$V_3>V_1>V_2$}
\opt{D}{$V_1>V_2>V_3$}
\end{qoptions}

\qn{9} Condenser A has a capacity of $15\mu F$ when it is filled with a medium of dielectric constant 15. Another condenser B has a capacity of $1\mu F$ with air between the plates. Both are charged separately by a battery of 100 V. After charging, both are connected in parallel without the battery and the dielectric medium being removed. The common potential now is :-
\begin{qoptions}
\opt{A}{400 V}
\opt{B}{800 V}
\opt{C}{1200 V}
\opt{D}{1600 V}
\end{qoptions}

\qn{10} A conducting wire bent in the form of a parabola $y^2=2x$ carries a current $i=2$ A as shown in figure. This wire is placed in a uniform magnetic field $\vec B=-4\hat k$ tesla. The magnetic force on the wire is (in newton) :-
\qfig{p1_phy_s1_q10_1}
\begin{qoptions}
\opt{A}{$-16\hat i$}
\opt{B}{$32\hat i$}
\opt{C}{$-32\hat i$}
\opt{D}{$16\hat i$}
\end{qoptions}

\qn{11} A wave in a string has an amplitude of 2cm. The wave travels in the +ve direction of x axis with a speed of 128 m/sec and it is noted that 5 complete waves fit in 4 m length of the string. The equation describing the wave is :-
\begin{qoptions}
\opt{A}{$y=(0.02)\ m\ \sin(7.85x-1005t)$}
\opt{B}{$y=(0.02)\ m\ \sin(7.85x+1005t)$}
\opt{C}{$y=(0.02)\ m\ \sin(15.7x-2010t)$}
\opt{D}{$y=(0.02)m\ \sin(15.7x+2010t)$}
\end{qoptions}

\qn{12} The number of possible natural oscillations of air column in a pipe closed at one end of length 85 cm whose frequencies lie below 1250 Hz are: (velocity of sound = 340 $ms^{-1}$)
\begin{qoptions}
\opt{A}{4}
\opt{B}{5}
\opt{C}{7}
\opt{D}{6}
\end{qoptions}

\qn{13} A thermodynamic process is shown in the figure. The pressures and volumes corresponding to some points in the figure are: $P_A=3\times10^4$ Pa, $P_B=8\times10^4$ Pa and $V_A=2\times10^{-3}\ m^3$, $V_D=5\times10^{-3}\ m^3$. In process AB, 600 J of heat is added to the system and in process BC, 200 J of heat is added to the system. The change in internal energy of the system in process AC would be:-
\qfig{p1_phy_s1_q13_1}
\begin{qoptions}
\opt{A}{560 J}
\opt{B}{800 J}
\opt{C}{600 J}
\opt{D}{640 J}
\end{qoptions}

\qn{14} The root mean square speed of the molecules of a gas is 1260 m/s. The most probable speed of the molecules is -
\begin{qoptions}
\opt{A}{1029 m/s}
\opt{B}{1161 m/s}
\opt{C}{1171 m/s}
\opt{D}{917 m/s}
\end{qoptions}

\qn{15} A heavy body of mass 25 kg is to be dragged along a horizontal plane ($\mu=1/\sqrt3$). The least force required is:
\begin{qoptions}
\opt{A}{25 kg f}
\opt{B}{2.5 kg f}
\opt{C}{12.5 kg f}
\opt{D}{$25/\sqrt3$ kg f}
\end{qoptions}

\qn{16} From a circle of radius a, an isosceles right angled triangle with the hypotenuse as the diameter of the circle is removed. The distance of the centre of gravity of the remaining portion from the centre of the circle is :-
\begin{qoptions}
\opt{A}{$\dfrac a3$}
\opt{B}{$\dfrac{(\pi-1)a}{6}$}
\opt{C}{$\dfrac{a}{3(\pi-1)}$}
\opt{D}{$\dfrac{a}{3(\pi+1)}$}
\end{qoptions}

\qn{17} A ball of mass m moving with a velocity v strikes the bob of a pendulum (mass m) at rest. If the collision is perfectly inelastic, the height to which both will rise is given by :-
\begin{qoptions}
\opt{A}{$\dfrac{v^2}{8g}$}
\opt{B}{$\dfrac{v^2}{4g}$}
\opt{C}{$\dfrac{v^2}{2g}$}
\opt{D}{$\dfrac{v^2}{g}$}
\end{qoptions}

\qn{18} Two discs have same mass and thickness. Their materials are of densities $\rho_1$ and $\rho_2$. The ratio of their moment of inertia about central axis will be-
\begin{qoptions}
\opt{A}{$\rho_1:\rho_2$}
\opt{B}{$\rho_1\rho_2:1$}
\opt{C}{$1:\rho_1\rho_2$}
\opt{D}{$\rho_2:\rho_1$}
\end{qoptions}

\qn{19} The earth moves around the sun in an elliptical orbit as shown in the figure. The ratio $\dfrac{OA}{OB}=x$. The ratio of the speeds of the earth at B and at A is :
\qfig{p1_phy_s1_q19_1}
\begin{qoptions}
\opt{A}{$\sqrt x$}
\opt{B}{$x$}
\opt{C}{$x^2$}
\opt{D}{$x\sqrt x$}
\end{qoptions}

\qn{20} A large open tank has two holes in the wall. One is a square hole of side L at a depth y from the top and the other is a circular hole of radius R at a depth 4y from the top. When the tank is completely filled with water, the quantities of water flowing out per second from both holes are the same. Then, R is equal to :-
\begin{qoptions}
\opt{A}{$\dfrac{L}{\sqrt{2\pi}}$}
\opt{B}{$2\pi L$}
\opt{C}{$L$}
\opt{D}{$\dfrac{L}{2\sqrt\pi}$}
\end{qoptions}

\sect{Section-II: Numerical Value}

\qn{1} A nuclear reactor generates power by fission of $_{92}U^{235}$ into two equal fragments of $_{46}Pd^{116}$ with the emission of two $\gamma$-rays of 5.2 MeV each and three neutrons. The average binding energies per nucleons of $_{92}U^{235}$ and $_{46}Pd^{116}$ are 7.2 MeV and 8.2 MeV respectively. The usable energy released per fission (in MeV) event is :-

\qn{2} Power lost in zener diode in the given circuit is $51.2\times10^{-x}$ w, find x :
\qfig{p1_phy_s2_q02_1}

\qn{3} In the given circuit as shown in figure the value of R(in $\Omega$) for which current through $4\Omega$ resistor is zero.
\qfig{p1_phy_s2_q03_1}

\qn{4} A conducting rod of 1 m length and 1 kg mass is suspended by two vertical wires through its ends. An external magnetic field of 2T is applied normal to the rod. Now, the current to be passed through the rod so as to make the tension in the wires zero is : (Take g = 10 $ms^{-2}$) :-

\qn{5} A narrow beam of ultrasonic waves reflects from a liver tumor at some depth inside the liver as shown in figure. The speed of the ultrasonic wave is 25\% less in the liver than in the surrounding medium. Then depth of the tumor in cm will be:
\qfig{p1_phy_s2_q05_1}

\qn{6} Two coherent sources of equal intensities produce a maximum of 100 units. If the amplitude of one of the sources is reduced by 20\%, then the maximum intensity produced will be :

\qn{7} A black body at a temperature of 1640 K has the wavelength corresponding to maximum emission equal to $1.75\mu m$. Assuming the moon to be a perfectly black body, the temperature of the moon, if the wavelength corresponding to maximum emission is $14.35\mu m$ is (in kelvin)

\qn{8} If the volume of a block of metal changes by 0.12\% when it is heated through $20^\circ C$, the coefficient of linear expansion (in $^\circ C^{-1}$) of the metal is $x\times10^{-5}$, find the value of x ?

\qn{9} The total weight of a piece of wood is 6 kg. In the floating state in water its 1/3 part remains inside the water. On this floating solid, what maximum weight is to be put such that the whole of the piece of wood is to be drowned in the water.

\qn{10} A 0.1 kg mass is suspended from a wire of negligible mass. The length of the wire is 1m and its cross-sectional area is $4.9\times10^{-7}\ m^2$. If the mass is pulled a little in the vertically downward direction and released, it performs simple harmonic motion of angular frequency 140 $rad\ s^{-1}$. If the Young's modulus of the material of the wire is $n\times10^9\ Nm^{-2}$, the value of n is

\partsec{2}{CHEMISTRY}

\sect{Section-I: Single Correct}

\qn{1} Gaseous $N_2O_4$ dissociates into gaseous $NO_2$ according to the reaction $N_2O_4(g)\rightleftharpoons2NO_2(g)$. At 300 K and 1 atm pressure, the degree of dissociation of $N_2O_4$ is 0.2. If one mole of $N_2O_4$ gas is contained in a vessel, then the density of the equilibrium mixture is :
\begin{qoptions}
\opt{A}{3.11 g/L}
\opt{B}{6.22 g/L}
\opt{C}{1.56 g/L}
\opt{D}{4.56 g/L}
\end{qoptions}

\qn{2} The emf in terms of oxidation potential of half cell given below is:\\
$Pt\,|\,H_2(g)\,|\,HCl$\\
\hspace*{1.2em}2 atm\hspace{2em}(0.02M)
\begin{qoptions}
\opt{A}{$+0.1093$ V}
\opt{B}{$-0.1143$ V}
\opt{C}{$+0.2286$ V}
\opt{D}{$-0.2286$ V}
\end{qoptions}

\qn{3} Electrolysis of hot aqueous solution of NaCl gives $NaClO_4$, i.e., sodium perchlorate, $NaCl+4H_2O\longrightarrow NaClO_4+4H_2\uparrow$. How many faraday are required to obtain 1000 g of sodium perchlorate :-
\begin{qoptions}
\opt{A}{65.3}
\opt{B}{40.3}
\opt{C}{18.3}
\opt{D}{31.6}
\end{qoptions}

\qn{4} If 20 mL of 0.1 M NaOH is added to 30 mL of 0.2 M $CH_3COOH$ ($pK_a=4.74$), the pH of the resulting solution is :
\begin{qoptions}
\opt{A}{4.44}
\opt{B}{9.56}
\opt{C}{8.96}
\opt{D}{9.26}
\end{qoptions}

\qn{5} If boiling point of an aqueous solution is $100.1^\circ C$, what is its freezing point? Given, enthalpy of fusion and vaporisation of water are 80 cal $g^{-1}$ and 540 cal $g^{-1}$ respectively :-
\begin{qoptions}
\opt{A}{$0.361^\circ C$}
\opt{B}{$-0.361^\circ C$}
\opt{C}{$-3.61^\circ C$}
\opt{D}{None of these}
\end{qoptions}

\qn{6} Complete combustion of 1.80 g of an oxygen containing compound $(C_xH_yO_z)$ gave 2.64 g of $CO_2$ and 1.08 g of $H_2O$. The percentage of oxygen in the organic compound is:
\begin{qoptions}
\opt{A}{51.63}
\opt{B}{63.53}
\opt{C}{53.33}
\opt{D}{50.33}
\end{qoptions}

\qn{7} $\Delta S_{Total}=-40$ kJ/mol-K\\
$\Delta H_{System}=2000$ kJ/mol\\
$T=400$ K\\
Find out value of $\Delta S_{System}$.
\begin{qoptions}
\opt{A}{$-35$ kJ/mol-K}
\opt{B}{$-10$ kJ/mol-K}
\opt{C}{$-40$ kJ/mol-K}
\opt{D}{$-5$ kJ/mol-K}
\end{qoptions}

\qn{8} Consider the following statement which of the following is not true
\begin{qoptions}
\opt{A}{$Eu^{+2}$ is a strong reducing agent (Z of Eu = 63)}
\opt{B}{$Ce^{+4}$ is a strong oxidising agent (Z of Ce = 58)}
\opt{C}{Curium have electronic configuration [Rn] $5f^76d^17s^2$ (Z of Cm = 96)}
\opt{D}{$Yb^{+2}$ is a oxidising agent (Z of Yb = 70)}
\end{qoptions}

\qn{9} Potassium permanganate acts as an oxidant in neutral, alkaline as well as acidic media. The final products obtained from it in the three conditions are, respectively :-
\begin{qoptions}
\opt{A}{$MnO_4^{2-}$, $Mn^{3+}$ and $Mn^{2+}$}
\opt{B}{$MnO_2$, $MnO_4^{2-}$ and $Mn^{2+}$}
\opt{C}{$MnO_2$, $MnO_4^{2-}$ and $Mn^{3+}$}
\opt{D}{$MnO_2$, $MnO_2$ and $Mn^{3+}$}
\end{qoptions}

\qn{10} Which of the following bonds has the highest energy
\begin{qoptions}
\opt{A}{Se--Se}
\opt{B}{Te--Te}
\opt{C}{S--S}
\opt{D}{O--O}
\end{qoptions}

\qn{11} The ionization potential of nitrogen is more than that of oxygen because of :-
\begin{qoptions}
\opt{A}{Greater attraction of electrons by the nucleus}
\opt{B}{Extra stability of the half filled p-orbitals}
\opt{C}{Smaller size of nitrogen}
\opt{D}{More penetration effect}
\end{qoptions}

\qn{12} The option with only amphoteric oxides is :-
\begin{qoptions}
\opt{A}{$Cr_2O_3$, NO, SnO, PbO}
\opt{B}{NO, $B_2O_3$, PbO, $SnO_2$}
\opt{C}{$Cr_2O_3$, BaO, SnO, $SnO_2$}
\opt{D}{ZnO, $Al_2O_3$, PbO, $PbO_2$}
\end{qoptions}

\qn{13} During change of $O_2$ to $O_2^-$ ion, the electron adds on which one of the following orbitals ?
\begin{qoptions}
\opt{A}{$\sigma^*$ orbital}
\opt{B}{$\sigma$ orbital}
\opt{C}{$\pi^*$ orbital}
\opt{D}{$\pi$ orbital}
\end{qoptions}

\qn{14} Which of the following complex is coloured and diamagnetic -
\begin{qoptions}
\opt{A}{$MnO_4^{2-}$}
\opt{B}{$[Ni(H_2O)_6]^{2+}$}
\opt{C}{$[CrCl_6]^{3-}$}
\opt{D}{$CrO_4^{2-}$}
\end{qoptions}

\qn{15} \qfig{p1_chem_s1_q15_1}
\begin{qoptions}
\opt{A}{\qfigi{p1_chem_s1_q15_2}}
\opt{B}{\qfigi{p1_chem_s1_q15_3}}
\opt{C}{\qfigi{p1_chem_s1_q15_5}}
\opt{D}{\qfigi{p1_chem_s1_q15_4}}
\end{qoptions}

\qn{16} Identify the major product in the following reaction.
\qfig{p1_chem_s1_q16_1}
\begin{qoptions}
\opt{A}{\qfigi{p1_chem_s1_q16_2}}
\opt{B}{\qfigi{p1_chem_s1_q16_3}}
\opt{C}{\qfigi{p1_chem_s1_q16_5}}
\opt{D}{\qfigi{p1_chem_s1_q16_4}}
\end{qoptions}

\qn{17} An organic compound A (Molecular formula $C_6H_{12}O$) does not change the colour of acidic dichromate solution. Compound A on treatment with $H_2SO_4$ produces alkene, which on oxidative ozonolysis gives a molecule ($C_6H_{10}O_3$) which gives positive iodoform test find a structure of A
\begin{qoptions}
\opt{A}{\qfigi{p1_chem_s1_q17_2}}
\opt{B}{\qfigi{p1_chem_s1_q17_1}}
\opt{C}{\qfigi{p1_chem_s1_q17_3}}
\opt{D}{\qfigi{p1_chem_s1_q17_4}}
\end{qoptions}

\qn{18} \qfig{p1_chem_s1_q18_1}
Structure of B will be :
\begin{qoptions}
\opt{A}{\qfigi{p1_chem_s1_q18_2}}
\opt{B}{\qfigi{p1_chem_s1_q18_3}}
\opt{C}{\qfigi{p1_chem_s1_q18_4}}
\opt{D}{\qfigi{p1_chem_s1_q18_5}}
\end{qoptions}

\qn{19} Salicin (structure given below) is a glycoside, found in the bark of willow tree, used in relieving pain. Observe the following reaction of salicin :
\qfig{p1_chem_s1_q19_1}
The incorrect statement is:
\begin{qoptions}
\opt{A}{P is D-glucose}
\opt{B}{Q is 2-hydroxybenzylalcohol}
\opt{C}{P will show mutarotation}
\opt{D}{P \& Q reduce tollens reagent}
\end{qoptions}

\qn{20} What could be the major product for the following reaction :-
\qfig{p1_chem_s1_q20_1}
\begin{qoptions}
\opt{A}{\qfigi{p1_chem_s1_q20_A}}
\opt{B}{\qfigi{p1_chem_s1_q20_3}}
\opt{C}{\qfigi{p1_chem_s1_q20_C}}
\opt{D}{\qfigi{p1_chem_s1_q20_5}}
\end{qoptions}

\sect{Section-II: Numerical Value}

\qn{1} For a first order reaction, time required for 99.0\% completion is `x' times for the time required for the completion of 90\% of the reaction then value of x is ?

\qn{2} 3g of acetic acid is added to 250 mL of 0.1 M HCl and the solution made up to 500 mL. To 20 mL of this solution $\dfrac12$ mL of 5 M NaOH is added. The pH of the solution is \_\_\_\_. (Nearest integer)\\
\relax[Given: $pK_a$ of acetic acid = 4.75, molar mass of acetic acid = 60 g/mol, $\log3=0.4771$]\\
Neglect any changes in volume

\qn{3} Given the following bond enthalpies :\\
$BE(N\equiv N)=942$ kJ/mol; $BE(H-H)=436$ kJ/mol;\\
$BE(N-N)=163$ kJ/mol; $BE(N-H)=390$ kJ/mol\\
Determine enthalpy change for the following polymerisation reaction per mole of $N_2(g)$ consumed -\\
$nN_2(g)+nH_2(g)\longrightarrow-(-NH-NH-)_n-$

\qn{4} Find $\Delta U$ for the process (in kJ) described by figure if heat supplied during the process is 50 kJ.
\qfig{p1_chem_s2_q04_1}

\qn{5} The value of $(n\times\ell\times m)$ for the last electron entered in K is

\qn{6} How many of the following ionic compound will have lattice energy less than $AlF_3$.\\
LiF, NaCl, $AlCl_3$, $Na_2O$, NaBr, KCl, LiI, $BeCl_2$, $MgF_2$

\qn{7} Number of $H_2O$ molecules attached as ligand in $CuSO_4\cdot5H_2O=x$\\
Number of $90^\circ$ bond angles in $XeF_4=y$\\
Number of ions in the aqueous solutions of $[Pt(NH_3)_6]Cl_4=z$ find $(x+y+z)=?$

\qn{8} How many Geometrical Isomers are possible for
\qfig{p1_chem_s2_q08_1}

\qn{9} How many moles of RMgX can be consumed by one mole of following compound when it reacts with excess of RMgX.
\qfig{p1_chem_s2_q09_1}

\qn{10} \qfig{p1_chem_s2_q10_1}
Number of possible alkenes formed in product mixture.

\partsec{3}{MATHEMATICS}

\sect{Section-I: Single Correct}

\qn{1} The value of $\cos\dfrac{2\pi}{7}+\cos\dfrac{4\pi}{7}+\cos\dfrac{6\pi}{7}+\cos\dfrac{7\pi}{7}$ is -
\begin{qoptions}
\opt{A}{1}
\opt{B}{$-1$}
\opt{C}{$1/2$}
\opt{D}{$-3/2$}
\end{qoptions}

\qn{2} If $0<\theta<\pi$, then minimum value of $3\sin\theta+\text{cosec}^3\theta$ is:-
\begin{qoptions}
\opt{A}{4}
\opt{B}{3}
\opt{C}{5}
\opt{D}{6}
\end{qoptions}

\qn{3} Mean deviation of first 7 natural no. about their A.M. is :-
\begin{qoptions}
\opt{A}{2}
\opt{B}{$\sqrt2$}
\opt{C}{$\dfrac{12}{7}$}
\opt{D}{0}
\end{qoptions}

\qn{4} Let $g(x)=f(x)+f(1-x)$ and $f''(x)>0$, $x\in(0,1)$. If g is decreasing in the interval $(0,\alpha)$ and increasing in the interval $(\alpha,1)$, then $\tan^{-1}(2\alpha)+\tan^{-1}\left(\dfrac1\alpha\right)+\tan^{-1}\left(\dfrac{\alpha+1}{\alpha}\right)$ is equal to :
\begin{qoptions}
\opt{A}{$\dfrac{3\pi}{2}$}
\opt{B}{$\pi$}
\opt{C}{$\dfrac{5\pi}{4}$}
\opt{D}{$\dfrac{3\pi}{4}$}
\end{qoptions}

\qn{5} Area of the region $\{(x,y):x^2+(y-2)^2\le4,\ x^2\ge2y\}$ is :
\begin{qoptions}
\opt{A}{$2\pi-\dfrac{16}{3}$}
\opt{B}{$\pi-\dfrac83$}
\opt{C}{$\pi+\dfrac83$}
\opt{D}{$2\pi+\dfrac{16}{3}$}
\end{qoptions}

\qn{6} Let $y=y(x)$ be the solution of the differential equations $\dfrac{dy}{dx}+\dfrac{5}{x(x^5+1)}y=\dfrac{(x^5+1)^2}{x^7}$, $x>0$. If $y(1)=2$, then $y(2)$ is equal to
\begin{qoptions}
\opt{A}{$\dfrac{637}{128}$}
\opt{B}{$\dfrac{679}{128}$}
\opt{C}{$\dfrac{693}{128}$}
\opt{D}{$\dfrac{697}{128}$}
\end{qoptions}

\qn{7} Let f and g be two functions defined by $f(x)=\begin{cases}x+1, & x<0\\|x-1|, & x\ge0\end{cases}$ and $g(x)=\begin{cases}x+1, & x<0\\1, & x\ge0\end{cases}$\\
Then $(g\circ f)(x)$ is
\begin{qoptions}
\opt{A}{Differentiable everywhere}
\opt{B}{Continuous everywhere but not differentiable exactly at one point}
\opt{C}{Not continuous at $x=-1$}
\opt{D}{Continuous everywhere but not differentiable at $x=1$}
\end{qoptions}

\qn{8} The domain of the function $f(x)=\dfrac{1}{\sqrt{[x]^2-3[x]-10}}$ is (where [x] denotes the greatest integer less than or equal to x)
\begin{qoptions}
\opt{A}{$(-\infty,-2)\cup(5,\infty)$}
\opt{B}{$(-\infty,-3]\cup[6,\infty)$}
\opt{C}{$(-\infty,-2)\cup[6,\infty)$}
\opt{D}{$(-\infty,-3]\cup(5,\infty)$}
\end{qoptions}

\qn{9} Assume that each born child is equally likely to be a boy or a girl. If two families have two children each, then the conditional probability that all children are girls given that at least two are girls is :
\begin{qoptions}
\opt{A}{$\dfrac{1}{11}$}
\opt{B}{$\dfrac{1}{17}$}
\opt{C}{$\dfrac{1}{10}$}
\opt{D}{$\dfrac{1}{12}$}
\end{qoptions}

\qn{10} Let $z_0$ be a root of the quadratic equation, $x^2+x+1=0$. If $z=3+6iz_0^{81}-3iz_0^{93}$, then $\arg z$ is equal to:
\begin{qoptions}
\opt{A}{$\dfrac\pi4$}
\opt{B}{$\dfrac\pi3$}
\opt{C}{0}
\opt{D}{$\dfrac\pi6$}
\end{qoptions}

\qn{11} If both the roots of the quadratic equation $x^2-mx+4=0$ are real and distinct and they lie in the interval $[1,5]$, then m lies in the interval:
\begin{qoptions}
\opt{A}{(4, 5)}
\opt{B}{(3, 4)}
\opt{C}{(5, 6)}
\opt{D}{$(-5,-4)$}
\end{qoptions}

\qn{12} The set of all values of $\lambda$ for which the system of linear equations.\\
$x-2y-2z=\lambda x$\\
$x+2y+z=\lambda y$\\
$-x-y=\lambda z$\\
has a non-trivial solution.
\begin{qoptions}
\opt{A}{contains more than two elements}
\opt{B}{is a singleton}
\opt{C}{is an empty set}
\opt{D}{contains exactly two elements}
\end{qoptions}

\qn{13} A group of students comprises of 5 boys and n girls. If the number of ways, in which a team of 3 students can randomly be selected from this group such that there is at least one boy and at least one girl in each team, is 1750, then n is equal to :
\begin{qoptions}
\opt{A}{25}
\opt{B}{28}
\opt{C}{27}
\opt{D}{24}
\end{qoptions}

\qn{14} Let $S_n$ denote the sum of the first n terms of an A.P. If $S_4=16$ and $S_6=-48$, then $S_{10}$ is equal to :
\begin{qoptions}
\opt{A}{$-200$}
\opt{B}{$-320$}
\opt{C}{400}
\opt{D}{$-400$}
\end{qoptions}

\qn{15} The remainder, when $7^{103}$ is divided by 17 is \_\_\_.
\begin{qoptions}
\opt{A}{12}
\opt{B}{13}
\opt{C}{7}
\opt{D}{5}
\end{qoptions}

\qn{16} If $\vec a=\hat i+2\hat k$, $\vec b=\hat i+\hat j+\hat k$, $\vec c=7\hat i-3\hat k+4\hat k$, $\vec r\times\vec b+\vec b\times\vec c=\vec0$ and $\vec r\cdot\vec a=0$ then $\vec r\cdot\vec c$ is equal to :
\begin{qoptions}
\opt{A}{34}
\opt{B}{12}
\opt{C}{36}
\opt{D}{30}
\end{qoptions}

\qn{17} Let $\vec a=4\hat i+3\hat j$ and $\vec b=3\hat i-4\hat j+5\hat k$ and $\vec c$ is a vector such that $\vec c\cdot(\vec a\times\vec b)+25=0$, $\vec c\cdot(\hat i+\hat j+\hat k)=4$ and projection of $\vec c$ on $\vec a$ is 1, then the projection of $\vec c$ on $\vec b$ equals :
\begin{qoptions}
\opt{A}{$\dfrac{5}{\sqrt2}$}
\opt{B}{$\dfrac15$}
\opt{C}{$\dfrac{1}{\sqrt2}$}
\opt{D}{$\dfrac{3}{\sqrt2}$}
\end{qoptions}

\qn{18} A straight line cuts off the intercepts $OA=a$ and $OB=b$ on the positive directions of x-axis and y-axis respectively. If the perpendicular from origin O to this line makes an angle of $\dfrac\pi6$ with positive direction of y-axis and the area of $\triangle OAB$ is $\dfrac{98}{3}\sqrt3$ then $a^2-b^2$ is equal to :
\begin{qoptions}
\opt{A}{$\dfrac{392}{3}$}
\opt{B}{196}
\opt{C}{$\dfrac{196}{3}$}
\opt{D}{98}
\end{qoptions}

\qn{19} The shortest distance between the lines $\dfrac{x-5}{1}=\dfrac{y-2}{2}=\dfrac{z-4}{-3}$ and $\dfrac{x+3}{1}=\dfrac{y+5}{4}=\dfrac{z-1}{-5}$ is :
\begin{qoptions}
\opt{A}{$7\sqrt3$}
\opt{B}{$5\sqrt3$}
\opt{C}{$6\sqrt3$}
\opt{D}{$4\sqrt3$}
\end{qoptions}

\qn{20} Let the ellipse $E:x^2+9y^2=9$ intersect the positive x- and y-axes at the points A and B respectively. Let the major axis of E be a diameter of the circle C. Let the line passing through A and B meet the circle C at the point P. If the area of the triangle which vertices A, P and the origin O is $\dfrac mn$, where m and n are coprime, then $m-n$ is equal to
\begin{qoptions}
\opt{A}{18}
\opt{B}{16}
\opt{C}{17}
\opt{D}{15}
\end{qoptions}

\sect{Section-II: Numerical Value}

\qn{1} If $\sec\theta+\tan\theta=\dfrac13$, then value of $16\cot^2\theta-5\sin\theta$ is :

\qn{2} If the domain of the function $f(x)=\sec^{-1}\left(\dfrac{2x}{5x+3}\right)$ is $[\alpha,\beta)\cup(\gamma,\delta]$ then $|3\alpha+10(\beta+\gamma)+21\delta|$ is equal to \_\_\_\_.

\qn{3} The number of relations, on the set $\{1,2,3\}$ containing (1,2) and (2,3), which are reflexive and transitive but not symmetric, is\_\_\_\_\_\_

\qn{4} If $\displaystyle\int_{0.15}^{0.15}\left|100x^2-1\right|dx=\dfrac{k}{3000}$, then k is equal to\_\_\_\_.

\qn{5} Let $I(x)=\displaystyle\int\sqrt{\dfrac{x+7}{x}}\,dx$ and $I(9)=12+7\log_e7$. If $I(1)=\alpha+7\log_e\left(1+2\sqrt2\right)$, then $\alpha^4$ is equal to\_\_\_\_\_\_\_.

\qn{6} Let [x] be the greatest integer $\le$ x. Then the number of points in the interval $(-2,1)$, where the function $f(x)=|[x]|+\sqrt{x-[x]}$ is discontinuous, is\_\_.

\qn{7} If some three consecutive in the binomial expansion of $(x+1)^n$ is powers of x are in the ratio $2:15:70$, then the average of these three coefficient is :-

\qn{8} If a and b are the roots of equation $x^2-7x-1=0$, then the value of $\dfrac{a^{21}+b^{21}+a^{17}+b^{17}}{a^{19}+b^{19}}$ is equal to \_\_.

\qn{9} If $P(h,k)$ be point on the parabola $x=4y^2$, which is nearest to the point $Q(0,33)$, then the distance of P from the directrix of the parabola $y^2=4(x+y)$ is equal to :

\qn{10} Let $P(a_1,b_1)$ and $Q(a_2,b_2)$ be two distinct points on a circle with center $C(\sqrt2,\sqrt3)$. Let O be the origin and OC be perpendicular to both CP and CQ. If the area of the triangle OCP is $\dfrac{\sqrt{35}}{2}$, then $a_1^2+a_2^2+b_1^2+b_2^2$ is equal to \_\_.

\end{multicols}
\end{document}
